% NeurIPS paper checklist -- answers and justifications only.
% Paste each Answer/Justification pair into the corresponding item of the template block,
% keeping the template's own question text and guidelines. Section references assume the
% reorganisation recommended in review_aug28: one Appendix A (Theory) and one Appendix B
% (Extended Experiments), with Appendix C merged away.

% ---------------------------------------------------------------- 1
\item {\bf Claims}
    \item[] Answer: \answerYes{}
    \item[] Justification: The abstract and introduction claim a population landing law and a
    flow-versus-do-nothing threshold for an informative source, a distinction between starting
    from the source and conditioning on it, and experiments testing both. These correspond to
    Theorem~1, Theorem~2, Propositions~2 and~3, and Sections~7.1--7.3. The abstract states that
    the real-data experiments test the qualitative picture rather than the numerical transfer of
    the synthetic law, which is the scope actually supported.

% ---------------------------------------------------------------- 2
\item {\bf Limitations}
    \item[] Answer: \answerYes{}
    \item[] Justification: Section~8 carries a dedicated Limitations paragraph. It states that
    the landing law and threshold are derived for the structured spherical model and that the
    exact factorization does not transfer to OpenCLIP embeddings; that curvature and von
    Mises--Fisher non-Gaussianity are not separated; that the regularized-source theory
    characterizes the degenerate limit rather than arbitrary finite regularization; that
    single-target cosine similarity can favour conditional averaging and is therefore not used as
    a standalone measure; and that the $1/t$ recovery singularity is specific to the interpolant
    analysed. The masked-diffusion result reported in the appendix, in which the
    source-conditioning penalty does not appear, is cited there as an explicit boundary on the
    claim.

% ---------------------------------------------------------------- 3
\item {\bf Theory assumptions and proofs}
    \item[] Answer: \answerYes{}
    \item[] Justification: Every formal statement is numbered and proved. Proposition~1 is proved
    in Appendix~A.2, Lemma~1 in Appendix~A.3, Theorem~1 in Appendix~A.4, Theorem~2 in
    Appendix~A.5, and Propositions~2 and~3 with Corollary~1 in Appendix~A.6. Assumptions are
    stated in the theorem statements: Proposition~1 requires a symmetric cross-covariance lying
    in the span of the two marginals; Theorem~1 is stated for the regularized-source family and
    assumes convergence of the terminal trajectories as the regularization vanishes; Theorem~2
    assumes $A\in(0,1)$ and $\jmath\in(0,\pi/2)$; Lemma~1 requires $d\ge4$. Proof sketches appear
    with each statement in the main text.

% ---------------------------------------------------------------- 4
\item {\bf Experimental result reproducibility}
    \item[] Answer: \answerYes{}
    \item[] Justification: Appendix~B specifies the generative law and its parameters, the
    operating point, the network architecture and width, the optimizer and learning rates, the
    parameterisation, batch size and step counts, the integration scheme and step count, the
    checkpoint-selection criterion with its separate evaluation seed, the number of seeds per
    configuration, and the evaluation protocol including the number of conditions and samples per
    condition. The real-data settings give the dataset splits and sizes, the OpenCLIP
    architecture and checkpoint, the blur construction, the degradation operator, and the
    pretrained image-to-image checkpoint. Code is released; see item~5.

% ---------------------------------------------------------------- 5
\item {\bf Open access to data and code}
    \item[] Answer: \answerYes{}
    \item[] Justification: An anonymized repository accompanies the submission containing the
    experiment scripts, the cluster job definitions, the analysis and plotting scripts that
    regenerate every figure and table from the saved outputs, and the saved result files
    themselves. All datasets used are public: CIFAR-10 and STL-10 are downloaded by the scripts,
    and the OpenCLIP and Stable Diffusion checkpoints are public. Each experiment is launched by
    a single documented command.

% ---------------------------------------------------------------- 6
\item {\bf Experimental setting/details}
    \item[] Answer: \answerYes{}
    \item[] Justification: The main text gives the operating point and the settings needed to
    read the results. Appendix~B gives the full details: training hyperparameters and the exact
    meaning of the training-budget comparison, the learning-rate and width grids, the integration
    scheme and step counts, the auxiliary-input constructions, the guidance formulation with its
    weights and the definition of the unguided reference, the source-informativeness
    parameterisation, and the dimensionality-estimation procedure with its sample sizes and
    eigenvalue cutoff.

% ---------------------------------------------------------------- 7
\item {\bf Experiment statistical significance}
    \item[] Answer: \answerYes{}
    \item[] Justification: The synthetic sweeps use five seeds per configuration, and the
    capacity table reports the gap in every learning-rate by width cell together with the
    smallest separation in standard errors of the mean across cells. The guidance table reports
    seed standard deviations for every entry. The super-resolution comparison uses three training
    seeds and states that the ordering holds in each seed separately. The pretrained
    image-to-image experiment reports standard errors obtained by bootstrapping the per-image
    paired difference over the evaluated images. The distributional measures are reported against
    a null floor obtained by splitting the target set in half, so each has its own zero. The
    evaluation protocol states the evaluation-noise floor of approximately $0.002$ on cosine
    quantities, measured by scoring one identical trained network twice; all reported gaps are at
    least an order of magnitude larger. Error bars are standard errors of the mean unless stated
    as standard deviations in the table caption.

% ---------------------------------------------------------------- 8
\item {\bf Experiments compute resources}
    \item[] Answer: \answerYes{}
    \item[] Justification: Appendix~B reports per-stage compute. The synthetic programme runs on
    a single GPU in approximately nine hours in total. Cluster experiments were run on NVIDIA
    A100 80GB nodes with PyTorch 2.1.0+cu121: the criterion-at-scale sweep, the three
    super-resolution resolutions (approximately two, two and four and a half hours), the
    six-factor sweep (approximately half an hour each), the three-learning-rate width sweep
    (approximately fifty minutes), the auxiliary-family retest, and the pretrained
    image-to-image experiment (approximately forty minutes). Total compute for the reported
    results is on the order of twenty-five GPU-hours. Preliminary and superseded runs, including
    an earlier single-metric version of the image-to-image experiment and an earlier
    super-resolution pass without the distributional measures, are not included in that figure
    and add roughly the same again.

% ---------------------------------------------------------------- 9
\item {\bf Code of ethics}
    \item[] Answer: \answerYes{}
    \item[] Justification: The research conforms to the NeurIPS Code of Ethics. It uses only
    public benchmark datasets and public pretrained checkpoints, involves no human subjects and
    no personally identifying data, and releases code rather than models. Anonymity is preserved
    in the submitted repository.

% ---------------------------------------------------------------- 10
\item {\bf Broader impacts}
    \item[] Answer: \answerYes{}
    \item[] Justification: The work is analytical and its practical consequence is a criterion
    for when \emph{not} to apply a generative transport, together with evidence that a pretrained
    image-to-image model can degrade an input it was applied to. The positive impact is reduced
    unnecessary generation in inverse problems, and a caution against evaluating restoration by a
    single-target similarity score, which the paper shows can reward averaged outputs. A negative
    consideration is that the same analysis identifies settings in which a generative model
    replaces rather than restores its input; the paper reports this as a failure mode to be
    avoided. The work introduces no new generative capability.

% ---------------------------------------------------------------- 11
\item {\bf Safeguards}
    \item[] Answer: \answerNA{}
    \item[] Justification: No models, datasets or scraped data are released. The released
    artefacts are analysis code and numerical result files for a controlled synthetic model and
    public benchmarks. The pretrained checkpoints used are already publicly available and are not
    redistributed.

% ---------------------------------------------------------------- 12
\item {\bf Licenses for existing assets}
    \item[] Answer: \answerYes{}
    \item[] Justification: CIFAR-10 and STL-10 are cited and used under their standard research
    terms. The OpenCLIP implementation is used under its MIT license and the ViT-B-32
    \texttt{laion2b\_s34b\_b79k} checkpoint under the terms of its release. The Stable Diffusion
    v1.5 checkpoint is used under the CreativeML OpenRAIL-M license. PyTorch, \texttt{diffusers},
    \texttt{transformers} and \texttt{torchvision} are used under their respective licenses.
    Versions are pinned in the repository. Each asset is cited at first use with its source URL.

% ---------------------------------------------------------------- 13
\item {\bf New assets}
    \item[] Answer: \answerYes{}
    \item[] Justification: The new asset is the code and the saved numerical outputs, released in
    the anonymized repository with a README documenting the command for each experiment, the
    structure of the result files, and the analysis scripts that regenerate each figure and
    table. No new dataset or model is introduced.

% ---------------------------------------------------------------- 14
\item {\bf Crowdsourcing and research with human subjects}
    \item[] Answer: \answerNA{}
    \item[] Justification: The work involves no crowdsourcing and no human subjects.

% ---------------------------------------------------------------- 15
\item {\bf Institutional review board (IRB) approvals}
    \item[] Answer: \answerNA{}
    \item[] Justification: The work involves no research with human subjects, so no IRB review
    was required.

% ---------------------------------------------------------------- 16
\item {\bf Declaration of LLM usage}
    \item[] Answer: \answerYes{}
    \item[] Justification: Large language models were used as a coding and analysis assistant
    throughout the experimental work: drafting and revising experiment and analysis scripts,
    inspecting saved outputs, and assisting with the exposition. They were also used to help
    derive and check several of the analytical results, each of which is stated as a numbered
    proposition with a written proof and is independently verified numerically against the same
    integrator used in the experiments; the verification scripts and their registers are included
    in the released repository. The research question, the experimental design, and the
    interpretation of the results are the authors'. No LLM is a component of the method under
    study.
